\documentclass[book.tex]{subfiles}

\begin{document}

\chapter{Lagrangian Mechanics}
\label{chap:lagrangian_mechanics}
\noindent\rule{11cm}{0.4pt} \\
\textbf{Prerequisites:} \autoref{chap:newtons_laws} \\
\textbf{Difficulty Level:} ** \\
\noindent\rule{11cm}{0.4pt}
\vspace{5mm}


Newtonian mechanics provides an extremely powerful framework to model mechanical systems, but has some limitations for more complex systems.
According to the Newtonian approach, the trajectory of a mass $m$ can be completely determined by solving the second law given an initial position and velocity. This framing may not be the most convenient to solve a complex physical system. Luckily, the Newtonian scheme is only one of many different formalism that can be used to solve a problem in classical mechanics. These different formulations are essentially the same theory, in the sense their domains and predictions are identical. Nevertheless, one or the other may be more effective for conceptual, computational or simply aesthetic reasons. 

In this chapter, we'll introduce the basics of the Lagrangian formalism, which provides an alternative mathematical representations of classical systems. In some cases, these alternative representations can be dramatically simpler than the Newtonian version. Looking ahead, Lagrangians and Hamiltonians (which we will introduce shortly) will prove useful for modeling more complex thermodynamical and quantum mechanical systems.

We may solve the mechanical problem considering the following approach: given the initial and final position of a mass, what distinguishes the real trajectory from all the other possible paths that connects these two points? We call this approach the Lagrangian formalism, and the answer to this question is based on the \textit{principle of stationary
action}.  Later in the book, we will see that in quantum mechanics, unlike classical mechanics, the evolution of a system is not restricted to a just single path, but rather to a number of trajectories, each one having a different probability.

\section{Lagrangian Mechanics}

One limitation of the Newtonian framework is that it requires explicitly tracking all the forces in a system. This can can get challenging to handle for complex systems, so we will often find it useful to swap to Lagrangian mechanics, which centers its model of reality around energy rather than forces. 




As a first guess, we define the type of the Lagrangian as a function 
\begin{align}
L: \mathcal{S} \times T \mathcal{S} \times \mathbb{R} \to \mathbb{R}
\end{align} 
That is, the Lagrangian is a function of the state space, derivative of the state space, and time to energy. In general we define the Lagrangian for the system as the difference between its kinetic energy, $T$, and potential energy, $V$
\begin{align}
L = T - V
\end{align}
where $T$ is the kinetic energy and $V$ is the potential energy. Using generalized coordinates for space, $q = (q_x,q_y,q_z,....)$, and time, $t$, we obtain

\begin{align}\label{lagrangian}
L(q, \dot{q}, t) &= T(q, \dot{q}, t) -V(q, t) \\
&= \frac{m}{2}\dot{q}^{2}-V(q,t)
\end{align}

The \textit{action}, $S$, is defined as the integral of the lagrangian between two given instants of time (where $\dot{q}=\frac{dq}{dt}$) as: 

\begin{align}\label{action}
S = \int_{t_1}^{t_2}\!L(\dot{q},q,t)\,dt
\end{align}

Note that the action is the integral of the Lagrangian over time. We see then that the action is a real valued quantity given $q$ and $\dot{q}$, so the action has a function type as follows
\begin{align}
    S: \mathcal{S} \times T\mathcal{S} \mapsto \mathbb{R}
\end{align} 

The principle of stationary (or least) action states that the path taken by the system between times $t_1$ and $t_2$, as shown in Fig.~\ref{fig:posa}, is the one for which the action is stationary (no change) to first order. Mathematically, for $\delta$ indicating a small change, this principle states:

\begin{align}\label{posaction}
\delta S = S[\bar{q}+\delta q]-S[\bar{q}] = 0
\end{align}

As the end points are fixed at $q_1$ and $q_2$, the perturbation has the condition $\delta q_1 = \delta q_2 = 0$. Using the definition of $S$ as in Eq.~\eqref{action}, we have:

\begin{align}
S[q+\delta q] &= \int_{t_1}^{t_2} L(\dot{q}+\delta \dot{q},q+\delta q,t)dt \\
&\approx \int_{t_1}^{t_2}L(\dot{q},q,t)+\delta \dot{q} \frac{\partial L}{\partial \dot{q}}+\delta q \frac{\partial L}{\partial q}dt \\
&= S[q] + \int_{t_1}^{t_2}\delta \dot{q} \frac{\partial L}{\partial \dot{q}}+\delta q \frac{\partial L}{\partial q}dt
\end{align}

\begin{marginfigure}%
% \includegraphics[width=\linewidth]{figures/Physics/mechanics/lagrangian_mechanics/posa.png}
\includegraphics[width=\linewidth]{figures/Physics/mechanics/lagrangian_mechanics/posa0.png}
\caption{Motion of a particle from $q_1$, at time $t_1$ to $q_2$, at time $t_2$ in the external potential $V(x,t)$. Among the several possible paths in $q$ and $t$ that the particle can traverse, the one denoted in red is the classical path that the particle chooses to traverse along, the other path curves (in purple) are not taken by the particle. }
\label{fig:posa}
\end{marginfigure}

Therefore, using integration by parts, the variation $\delta S$ can be written as:

\begin{align}\label{posaction2}
\delta S &= \int_{t_1}^{t_2}\!\delta \dot{q} \frac{\partial L}{\partial \dot{q}}+\delta q \frac{\partial L}{\partial q}\,dt = \delta q \left. \frac{\partial L}{\partial \dot{q}} \right|_{t_1}^{t_2}  - \int_{t_1}^{t_2}\!\delta q \left[\frac{d}{dt} \frac{\partial L}{\partial \dot{q}}- \frac{\partial L}{\partial q}\right]\,dt
\end{align}

The first term in Eq.~\ref{posaction2} is zero as $\delta q_1 = \delta q_2 = 0$. Therefore, regardless of $\delta q$, the path with the minimum action will satisfy the condition:

\begin{align}\label{zeroaction}
\frac{d}{dt} \frac{\partial L}{\partial \dot{q}} - \frac{\partial L}{\partial q} = 0
\end{align}

This equation is known as Lagrange's equations or the Euler-Lagrange equations. Note that above we treat $q$ and $\dot{q}$ as independent arguments to $L$. Finally, using the definition of $L$, we have the equation of motion as:

\begin{align}\label{motion}
m \frac{d^2q}{dt^2} + \frac{\partial V}{\partial q}= 0
\end{align}

Defining the force due to the external potential to be $\vec{F} =- \frac{\partial V}{\partial q}$ we have $- \frac{\partial V}{\partial q} = m\vec{a}$, which is Newton's second law of motion. Now, we consider example problems to draw some conclusions about classical mechanics.


\subsection{A free particle}

Consider a free particle with 1-D motion along the $x$ axis and the external potential $V(x,t)=0$. In this case, the state space is $\mathcal{S} : \mathbb{R}$. The Lagrangian is then given by
\begin{align}
L(x, \dot{x}) &= \frac{1}{2} m \dot{x}^2
\end{align}

Therefore, its equation of motion will be:

\begin{align}\label{exmp1}
\frac{d}{dt} \frac{\partial L}{\partial \dot{x}} - \frac{\partial L}{\partial x} &= 0 \\
m \frac{d^2x}{dt^2} &= 0 \\ x(t) &= C_1t + C_2 
\end{align}

Where $C_1$, $C_2$ are constants determined using the initial conditions. Considering $x=0$ at $t=0$ and $\dot{x}=v=0$ at $t=0$, we get $x=vt$ as the equation describing the particle's motion. This result is in agreement with Newton's law of motion.

\subsection{A particle in a harmonic potential field}

Consider the same particle as in the previous section, but with the external potential $V(x,t)= \frac{k}{2}x^2$. The Lagrangian is now given by
\begin{align}
    L(x, \dot{x}) &= \frac{1}{2}m\dot{x}^2 - \frac{k}{2}x^2
\end{align}

This will result in an equation of motion given by:

\begin{align}
\frac{d}{dt}\frac{\partial L}{\partial \dot{x}} - \frac{\partial L}{\partial x} &= 0 \\
m \frac{d^2x}{dt^2} + kx &= 0 \\ \frac{d^2x}{dt^2} + \nu^2x &= 0 
\end{align}

Where $\nu = \sqrt{\frac{k}{m}}$ is called the characteristic frequency. Let's quickly consider the units of $k$ to verify that it has the dimensions for a frequency::
\begin{lstlisting}
    [$\nu$] = $\sqrt{\frac{\textrm{[k]}}{\textrm{[m]}}}$ = $\sqrt{\frac{\textrm{[M]}}{\textrm{[T]}^2 \textrm{[M]}}}$ = $\sqrt{\frac{1}{\textrm{[T]}^2}}$ = [T]$^{-1}$
\end{lstlisting}

This is the same equation as we saw earlier. Considering $x=0$ at $t=0$ and $\dot{x}=v=v_0$ at $t=0$, we get the equations describing the particle's motion as:

\begin{align}
x(t) &= A\sin(\nu t)+B\cos(\nu t) = \frac{v_0}{\nu}\sin(\nu t) \\
v(t) &= \frac{dx}{dt} = v_0\cos(\nu t) 
\end{align}


\subsection{A box sliding down a ramp}

% TODO: This figure is sourced from \url{http://www.dzre.com/alex/P441/lectures/lec_18.pdf}, Example 1. Replace with our own illustration.
\begin{figure}
   % \includegraphics[width=\linewidth]{figures/Physics/mechanics/lagrangian_mechanics/box_on_ramp.png}
   \centering\includegraphics[width=0.5\linewidth]{figures/Physics/mechanics/lagrangian_mechanics/Frictionless Box.png}
   \caption{A frictionless box on a ramp.}
\label{fig:box_on_ramp}
\end{figure}

Let's now introduce a physical theory for a simple system with a box moving down a ramp. This simple model has no friction, but both the ramp and box in question can move. The state space $\mathcal{S}$ for this system is given by
\begin{align}
    \mathcal{S} = \mathbb{R} \times \mathbb{R} 
\end{align}
Each state consists of a tuples $(x_1, x_2)$ of positions for the ramp and the box. Note that for this system, position is one dimensional and not three dimensional since there's only 1 axis of motion for both box and ramp. Note this will constrain the domain of applicability for this physical theory, but will give us the nice side effect of simplicity. 

Let $v_1$ be the velocity of the ramp, and $v_2$ the velocity of the box relative to the ramp. The energy is given by
\begin{align}
E = \frac{1}{2}Mv_1^2 + \frac{1}{2}mv_2^2 \sin^2 \theta + \frac{1}{2} m(v_2 \cos \theta + v_1)^2 + mg(H - x_2 \sin \theta)
\label{ramp_energy}
\end{align}
Here $\theta$ is the angle of the ramp. In the Lagrangian formulation, the time evolution equations are derived from the Lagrangian.
\begin{align}
\mathcal{L}(x_1, x_2, v_1, v_2, t) &= T - V\\
&= \frac{1}{2}Mv_1^2 + \frac{1}{2}mv_2^2 \sin^2 \theta + \frac{1}{2} m(v_2 \cos \theta + v_1)^2  - mg(H - x_2 \sin \theta)
\end{align}
With this Lagrangian, the evolution is governed by the solutions to the Euler-Lagrange equations
\begin{align}
    \frac{d}{dt} \frac{\partial \mathcal{L}}{\partial v_1} &= \frac{\partial \mathcal{L}}{\partial x_1} \\
    \frac{d}{dt} \left ( Mv_1 + m(v_2\cos \theta +v_1) \right ) &= 0 \\
    M a_1 + ma_2 \cos \theta + m a_1 &= 0 \\
    (M + m) a_1 + m \cos \theta a_2 &= 0
\label{ramp:euler:lagrange}
\end{align}

\begin{align}
\frac{d}{dt} \frac{\partial \mathcal{L}}{\partial v_2} &= \frac{\partial L}{\partial x_2} \\
\frac{d}{dt} \left ( m v_2 \sin^2 \theta + m(v_2 \cos \theta + v_1) \right ) &= mg\sin \theta \\
m a_2 \sin^2 \theta + m a_2 \cos \theta + m a_1  &= mg \sin \theta \\
m a_1 + (m \sin^2 \theta + m \cos \theta) a_2 &= mg \sin \theta
\end{align}
We consequently have a pair of coupled nonlinear differential equations. We can solve these equations using the techniques we have learned about previously.

For the domain of applicability, we must have 
\begin{align}
x_2 < \frac{H}{\sin \theta}
\end{align}
since the theory is only valid when the box is on the ramp. We also assume that $\|v_1\|, \|v_2\| < C$ where $C$ is some constant. As usual, this is done to avoid relativistic effects.
Parameters are given by
\begin{align}\mathcal{W} = \{M, m, \theta, H, g, \vec{x_1}(0), \vec{x_2}(0)\}\end{align}
where $m$ is the mass of the box, $M$ the mass of the ramp, $\theta$ the angle of inclination of the ramp, $H$ the height of the ramp, and $g$ is the gravitational constant on earth. $\vec{x_1}(0)$ and $\vec{x_2}(0)$ are the initial positions of the box and ramp.
Hyperparameters are given by
\begin{align}\mathcal{H} = \{C\}\end{align}

How can we solve this physical theory? In order to solve this theory, we need to solve Equations \ref{ramp:euler:lagrange} to find equations for $x_1(t)$ and $x_2(t)$. We leave this solution to the exercises.

\subsection{Exercises}
\begin{enumerate}

\item Derive the equation for energy in Equation \ref{ramp_energy}.
\item Solve the Equations \ref{ramp:euler:lagrange} for $x_1(t)$ and $x_2(t)$. Hint: Derive the following equations for the acceleration by solving the Euler-Lagrange equations.
\begin{align}
    a_1(t) &= \frac{-g \sin \theta \cos \theta}{\frac{m+M}{m} - \cos^2(\theta)} \\
    a_2(t) &= \frac{g \sin \theta}{1 - \frac{m\cos^2(\theta)}{m+M}}
\end{align}
\end{enumerate}

\end{document}